\documentclass[11pt]{article}
\usepackage{amsmath,amssymb,amsthm}
\newtheorem{proposition}{Proposition}
\begin{document}
\title{A correction to the open gadget assertion}
\author{}
\date{}
\maketitle
This fragment is separate from the preserved source of
\emph{Rigidity of harmonic functions on the supercritical percolation cluster},
arXiv:2303.04736v2 (29 March 2024). It corrects only the gadget proposition
at source lines 1141--1144 and the direction of divisibility at line 1319.
The proposition as printed on the closed rectangle is false: values at unused
closed sites can make a function nonconstant without changing either terminal value.
The original source and its abstract are not revised by this fragment.

For integers $d\geq2$ and $n\geq1$, put
\[
V_n=\bigl(([1,4]\cap\mathbb Z)\times\{1\}\bigr)
 \cup\bigl(\{2,3\}\times([1,n]\cap\mathbb Z)\bigr),
\]
and regard this set as a subset of $\mathbb Z^d$ by setting every coordinate
after the second equal to zero. Set $s=(1,1,0,\ldots,0)$ and
$t=(4,1,0,\ldots,0)$. Join precisely those pairs in $V_n$ whose
$\ell^1$ distance in $\mathbb Z^d$ is one. For a real function on
$\mathbb Z^d$, the induced graph Laplacian is
\[
\Delta_{V_n}u(x)=\sum_{\substack{y\in V_n\\|y-x|_1=1}}(u(y)-u(x)),
\qquad x\in V_n.
\]
Thus the neighbors are taken in the actual open carrier, including the top rung.
The two terminals are the only vertices at which no harmonic equation is required.

\begin{proposition}
There is one constant $C>0$ such that, for every $d\geq2$, every integer
$n\geq1$, and every $u : \mathbb Z^d\to\mathbb R$ satisfying
$u(x)\in\mathbb Z$ for all $x\in V_n$,
$\Delta_{V_n}u(x)=0$ for every $x\in V_n\setminus\{s,t\}$,
and $u(x)\ne u(y)$ for some $x,y\in V_n$, one has
\[
 C3^n\leq |u(t)-u(s)|.
\]
In fact $C=1$ works, and $|u(t)-u(s)|\geq A_{n+1}$, where
$A_0=A_1=1$ and $A_{k+2}=4A_{k+1}-A_k$ for every integer $k\geq0$.
\end{proposition}

\begin{proof}
Write $\delta_j=u(3,j,0,\ldots,0)-u(2,j,0,\ldots,0)$ for
$1\leq j\leq n$. If $n=1$, the harmonic equations at $(2,1)$ and
$(3,1)$ give $u(t)-u(s)=3\delta_1=A_2\delta_1$.
If $n\geq2$, subtracting the two column equations at the top rung,
at an interior rung, and at the bottom rung, respectively, gives
\[
\delta_{n-1}=3\delta_n,\qquad
\delta_{j-1}=4\delta_j-\delta_{j+1}\ (1<j<n),\qquad
u(t)-u(s)=4\delta_1-\delta_2.
\]
Backward induction with the defining recurrence of $A$ yields
$\delta_j=A_{n-j+1}\delta_n$ for every $1\leq j\leq n$.
Consequently, in both height cases,
\[
 u(t)-u(s)=A_{n+1}\delta_n.
\]
The finite graph is connected, since each column is joined by vertical edges,
each rung joins the two columns, and the two leaves attach at the bottom.
If the terminals had the same value, the maximum principle would force the
function to be constant: a maximum or minimum different from that terminal
value occurs at a nonterminal, and the harmonic equation propagates that
extreme along every adjacent edge until a terminal is reached.
Nonconstancy therefore implies $u(t)\ne u(s)$, hence $\delta_n\ne0$.
The integer values on $V_n$ give $|\delta_n|\geq1$.
Finally, the recurrence gives $A_{k+1}\geq3^k$ for every $k\geq0$:
$A_1=1$, $A_2=3$, and positivity and
$A_{k+2}\geq3A_{k+1}$ propagate since $A_k\leq A_{k+1}$.
This proves the asserted stronger bound and the uniform choice $C=1$.
\end{proof}

At source line 1319, replace the claim that the normalized value $u(t)$
divides $A_{n+1}$ by the claim that $A_{n+1}$ divides $u(t)$ after
subtracting $u(s)$. The displayed terminal transfer identity proves this
direction directly. Integer recovery from real pointwise integer values
introduces no new premise: each value has a unique integer representative,
by injectivity of the cast from $\mathbb Z$ to $\mathbb R$.
No cluster attachment, isolation of nonterminal bonds, restriction of
cluster harmonicity, or cylinder probability is asserted here. Those
applications and the three main rigidity theorems are outside this fragment.
\end{document}
